% Generated draft: requires mathematical and source review.
\documentclass[11pt]{article}
\usepackage[margin=25mm]{geometry}
\usepackage{fontspec}
\setmainfont{Latin Modern Roman}
\usepackage{amsmath,amssymb,mathtools,unicode-math}
\setmathfont{Latin Modern Math}
\usepackage{enumitem,needspace,xurl,hyperref}
\setlength{\parindent}{0pt}
\setlength{\parskip}{5pt}
\setlength{\emergencystretch}{3em}
\newcommand{\N}{\mathbb{N}}
\newcommand{\Z}{\mathbb{Z}}
\newcommand{\Q}{\mathbb{Q}}
\newcommand{\R}{\mathbb{R}}
\newcommand{\C}{\mathbb{C}}
\newcommand{\F}{\mathbb{F}}
\newcommand{\A}{\mathbb{A}}
\newcommand{\PP}{\mathbb{P}}
\newcommand{\E}{\mathbb{E}}
\newcommand{\eps}{\varepsilon}
\newcommand{\dd}{\,\mathrm{d}}
\begin{document}
\section*{UnsolvedMath ID: 30001070}
\textbf{Source problem code:} catalogue entry
\subsection{Definitions and mathematical statement}
For an integer \(n\geq1\), let \(B^n=\{x\in\R^n:\|x\|_2\leq1\}\). Let \(P\subseteq\R^n\) be a full-dimensional convex polytope, hence bounded, with exactly \(2n\) facets and with \(B^n\subseteq P\). Define its outer radius about the fixed ball center by
\[
 R_0(P)=\max_{x\in P}\|x\|_2.
\]
This maximum is attained at a vertex. It is not an optimization over possible centers of enclosing balls.

The conjecture asserts
\[
 \forall n\geq1\ \forall P\supseteq B^n\text{ with exactly }2n\text{ facets},
 \qquad R_0(P)\geq\sqrt n,
\]
with the full equality characterization
\[
 R_0(P)=\sqrt n\quad\Longleftrightarrow\quad
 P=Q[-1,1]^n\text{ for some }Q\in O(n).
\]
The right side describes a Euclidean cube centered at the origin, of side length two, circumscribed about the unit ball; all its facets are tangent to \(B^n\). A merely combinatorial cube or general parallelotope is not the stated equality case. The general hypothesis only requires containment of the ball, not tangency of every facet or central symmetry of \(P\).

Equivalently, with unit outward facet normals one may write
\[
 P=\bigcap_{i=1}^{2n}\{x:\langle a_i,x\rangle\leq b_i\},
 \qquad\|a_i\|_2=1,\quad b_i\geq1,
\]
with a genuine facet for each inequality; the \(b_i\) are not assumed all equal to one.

The original explicitly says the conjecture is open for \(n\geq5\). The supplied assessment dated 2026-08-21 instead refers to a blocking-number conjecture as open even in dimension three and reports no later proof in the stated generality. These dimensional descriptions are not reconciled in the supplied record and may concern different formulations. They are preserved as source qualifications, without asserting a verified status in dimensions three, four, or higher.
\subsection{Short English statement}
Any polytope with exactly twice as many facets as the dimension and containing the unit ball is conjectured to reach distance at least \(\sqrt n\) from the origin, with equality only for the tangent Euclidean cube. The source's original and dated assessment disagree about the dimensional scope of the open cases.
\subsection{Sources}
\begin{itemize}
\item Local UnsolvedMath registry, permanent ID \(30001070\), exact source code \texttt{catalogue entry}, proposed year 2008; no individual proposer is supplied.
\item \emph{Discrete Geometry (2009)}, Oberwolfach Reports, DOI 10.4171/owr/2008/44, pp. 2477--2556. \url{https://doi.org/10.4171/owr/2008/44}.
\item The supplied \texttt{corrected\_verified} note reports meaning-preserving semantic and notation cleanup. The original says open for \(n\geq5\); the open label and assessment checked 2026-08-21 describe a blocking-number question open even in dimension three. This is supplied dated curation, not independently verified, and the discrepancy requires source review.
\end{itemize}
\end{document}
